% cnreport-preamble.tex — 中文技术报告导言区(Cosmos/arXiv 单栏风格)
% 用法: \documentclass[11pt]{ctexart} 之后 % cnreport-preamble.tex — 中文技术报告导言区(Cosmos/arXiv 单栏风格)
% 用法: \documentclass[11pt]{ctexart} 之后 % cnreport-preamble.tex — 中文技术报告导言区(Cosmos/arXiv 单栏风格)
% 用法: \documentclass[11pt]{ctexart} 之后 % cnreport-preamble.tex — 中文技术报告导言区(Cosmos/arXiv 单栏风格)
% 用法: \documentclass[11pt]{ctexart} 之后 \input{cnreport-preamble}
% 编译: tectonic(XeTeX 引擎),需要系统安装 Noto Serif/Sans CJK SC 字体
\usepackage[a4paper,top=2.7cm,bottom=2.7cm,left=2.5cm,right=2.5cm,headheight=15pt]{geometry}
\usepackage{fontspec}
% 正文:拉丁衬线 + 思源宋体;标题:拉丁无衬线 + 思源黑体
% 按文件名从 TeX 字体树加载(tectonic bundle 自带 TeX Gyre;系统字体需 Noto CJK SC)
\setmainfont{texgyretermes}[Extension=.otf,UprightFont=*-regular,BoldFont=*-bold,ItalicFont=*-italic,BoldItalicFont=*-bolditalic]
\setsansfont{texgyreheros}[Extension=.otf,UprightFont=*-regular,BoldFont=*-bold,ItalicFont=*-italic,BoldItalicFont=*-bolditalic]
\setmonofont{texgyrecursor}[Extension=.otf,UprightFont=*-regular,BoldFont=*-bold,Scale=0.88]
\setCJKmainfont{Noto Serif CJK SC}[BoldFont={Noto Serif CJK SC Bold}]
\setCJKsansfont{Noto Sans CJK SC}[BoldFont={Noto Sans CJK SC Bold}]
\setCJKmonofont{Noto Sans Mono CJK SC}
% 带圈数字 ①-⑳ 等划入 CJK 类,用 CJK 字体渲染(xeCJK 默认按拉丁处理)
\xeCJKDeclareCharClass{CJK}{"2460 -> "24FF}

\usepackage{amsmath,amssymb,bm}
\usepackage{booktabs}
\usepackage{graphicx}
\usepackage{xcolor}
\usepackage{titlesec}
\usepackage{fancyhdr}
\usepackage{enumitem}
\usepackage{caption}
\usepackage{listings}

% ---- 配色(主色可改) ----
\definecolor{accent}{HTML}{0E5FA8}     % 强调蓝(链接/节号)
\definecolor{good}{HTML}{0E7A4D}       % 正向结果
\definecolor{warn}{HTML}{B45309}       % 负向/警示
\definecolor{faint}{HTML}{6B7889}      % 弱化文字
\definecolor{codebg}{HTML}{F0F2F5}
\definecolor{rule}{HTML}{C9CFD8}

\usepackage[colorlinks=true,linkcolor=accent,citecolor=accent,urlcolor=accent]{hyperref}

% ---- 章节标题:无衬线粗体,节号着色 ----
\titleformat{\section}{\Large\bfseries\sffamily}{\textcolor{accent}{\thesection}}{0.6em}{}
\titleformat{\subsection}{\large\bfseries\sffamily}{\textcolor{accent}{\thesubsection}}{0.5em}{}
\titleformat{\subsubsection}{\normalsize\bfseries\sffamily}{\textcolor{accent}{\thesubsubsection}}{0.5em}{}
\titlespacing{\section}{0pt}{1.6em}{0.7em}
\titlespacing{\subsection}{0pt}{1.2em}{0.5em}

% ---- 页眉页脚(Cosmos 式:左标题右日期,下细线) ----
\newcommand{\runningtitle}{}  % 主文档用 \renewcommand 设置
\newcommand{\reportdate}{}
\newcommand{\reportfooter}{}
\fancypagestyle{cnreport}{%
  \fancyhf{}
  \fancyhead[L]{\small\sffamily \runningtitle}
  \fancyhead[R]{\small\sffamily \reportdate}
  \fancyfoot[L]{\small\color{faint}\sffamily \reportfooter}
  \fancyfoot[R]{\small\sffamily \thepage}
  \renewcommand{\headrulewidth}{0.4pt}
  \renewcommand{\footrulewidth}{0pt}
}
\pagestyle{cnreport}

% ---- 题注与列表 ----
\captionsetup{font=small,labelfont={bf,sf},skip=6pt}
\setlist{itemsep=2pt,topsep=4pt,parsep=0pt}

% ---- 代码 ----
\lstset{
  basicstyle=\ttfamily\small,
  backgroundcolor=\color{codebg},
  frame=single, framerule=0.3pt, rulecolor=\color{rule},
  xleftmargin=6pt, xrightmargin=6pt,
  breaklines=true, columns=fullflexible, keepspaces=true,
  aboveskip=8pt, belowskip=8pt,
}

% ---- 报告头(kicker / 标题 / 副题 / 署名 / 摘要 / 关键词) ----
% 用法: \reporthead{kicker-left}{kicker-right}{标题}{副题}{署名行}
\newcommand{\reporthead}[5]{%
  \begin{center}
  {\small\sffamily\bfseries\color{accent} #1 \hfill #2}\\[2pt]
  {\color{black}\rule{\linewidth}{1.2pt}}\\[14pt]
  {\LARGE\sffamily\bfseries #3\par}\vspace{8pt}
  {\large\color{faint} #4\par}\vspace{10pt}
  {\small\sffamily\color{faint} #5\par}\vspace{4pt}
  {\color{rule}\rule{\linewidth}{0.4pt}}
  \end{center}\vspace{6pt}
}
% 摘要块: \begin{cnabstract} ... \end{cnabstract}
\newenvironment{cnabstract}{%
  \begin{quote}\small\color{faint}\sffamily\bfseries\color{black} 摘要 \quad\normalfont\rmfamily\color{faint!80!black}%
}{\end{quote}\vspace{2pt}}
% 重点结论框: \begin{quotebox} ... \end{quotebox}
\newenvironment{quotebox}{%
  \begin{quote}\small
  {\color{rule}\rule{\linewidth}{0.8pt}}\par\vspace{2pt}%
}{\par\vspace{2pt}{\color{rule}\rule{\linewidth}{0.8pt}}\end{quote}\vspace{2pt}}

\newcommand{\keywords}[1]{\begin{quote}\small\color{faint}{\bfseries\sffamily\color{black} 关键词:\ } #1\end{quote}}

% 表内彩色标记
\newcommand{\upgood}[1]{\textcolor{good}{\bfseries #1}}
\newcommand{\downbad}[1]{\textcolor{warn}{\bfseries #1}}
\newcommand{\dimtext}[1]{\textcolor{faint}{#1}}

\setlength{\parskip}{2pt}
\linespread{1.12}
\emergencystretch=3em  % CJK 段落容忍度,消除 overfull

% 编译: tectonic(XeTeX 引擎),需要系统安装 Noto Serif/Sans CJK SC 字体
\usepackage[a4paper,top=2.7cm,bottom=2.7cm,left=2.5cm,right=2.5cm,headheight=15pt]{geometry}
\usepackage{fontspec}
% 正文:拉丁衬线 + 思源宋体;标题:拉丁无衬线 + 思源黑体
% 按文件名从 TeX 字体树加载(tectonic bundle 自带 TeX Gyre;系统字体需 Noto CJK SC)
\setmainfont{texgyretermes}[Extension=.otf,UprightFont=*-regular,BoldFont=*-bold,ItalicFont=*-italic,BoldItalicFont=*-bolditalic]
\setsansfont{texgyreheros}[Extension=.otf,UprightFont=*-regular,BoldFont=*-bold,ItalicFont=*-italic,BoldItalicFont=*-bolditalic]
\setmonofont{texgyrecursor}[Extension=.otf,UprightFont=*-regular,BoldFont=*-bold,Scale=0.88]
\setCJKmainfont{Noto Serif CJK SC}[BoldFont={Noto Serif CJK SC Bold}]
\setCJKsansfont{Noto Sans CJK SC}[BoldFont={Noto Sans CJK SC Bold}]
\setCJKmonofont{Noto Sans Mono CJK SC}
% 带圈数字 ①-⑳ 等划入 CJK 类,用 CJK 字体渲染(xeCJK 默认按拉丁处理)
\xeCJKDeclareCharClass{CJK}{"2460 -> "24FF}

\usepackage{amsmath,amssymb,bm}
\usepackage{booktabs}
\usepackage{graphicx}
\usepackage{xcolor}
\usepackage{titlesec}
\usepackage{fancyhdr}
\usepackage{enumitem}
\usepackage{caption}
\usepackage{listings}

% ---- 配色(主色可改) ----
\definecolor{accent}{HTML}{0E5FA8}     % 强调蓝(链接/节号)
\definecolor{good}{HTML}{0E7A4D}       % 正向结果
\definecolor{warn}{HTML}{B45309}       % 负向/警示
\definecolor{faint}{HTML}{6B7889}      % 弱化文字
\definecolor{codebg}{HTML}{F0F2F5}
\definecolor{rule}{HTML}{C9CFD8}

\usepackage[colorlinks=true,linkcolor=accent,citecolor=accent,urlcolor=accent]{hyperref}

% ---- 章节标题:无衬线粗体,节号着色 ----
\titleformat{\section}{\Large\bfseries\sffamily}{\textcolor{accent}{\thesection}}{0.6em}{}
\titleformat{\subsection}{\large\bfseries\sffamily}{\textcolor{accent}{\thesubsection}}{0.5em}{}
\titleformat{\subsubsection}{\normalsize\bfseries\sffamily}{\textcolor{accent}{\thesubsubsection}}{0.5em}{}
\titlespacing{\section}{0pt}{1.6em}{0.7em}
\titlespacing{\subsection}{0pt}{1.2em}{0.5em}

% ---- 页眉页脚(Cosmos 式:左标题右日期,下细线) ----
\newcommand{\runningtitle}{}  % 主文档用 \renewcommand 设置
\newcommand{\reportdate}{}
\newcommand{\reportfooter}{}
\fancypagestyle{cnreport}{%
  \fancyhf{}
  \fancyhead[L]{\small\sffamily \runningtitle}
  \fancyhead[R]{\small\sffamily \reportdate}
  \fancyfoot[L]{\small\color{faint}\sffamily \reportfooter}
  \fancyfoot[R]{\small\sffamily \thepage}
  \renewcommand{\headrulewidth}{0.4pt}
  \renewcommand{\footrulewidth}{0pt}
}
\pagestyle{cnreport}

% ---- 题注与列表 ----
\captionsetup{font=small,labelfont={bf,sf},skip=6pt}
\setlist{itemsep=2pt,topsep=4pt,parsep=0pt}

% ---- 代码 ----
\lstset{
  basicstyle=\ttfamily\small,
  backgroundcolor=\color{codebg},
  frame=single, framerule=0.3pt, rulecolor=\color{rule},
  xleftmargin=6pt, xrightmargin=6pt,
  breaklines=true, columns=fullflexible, keepspaces=true,
  aboveskip=8pt, belowskip=8pt,
}

% ---- 报告头(kicker / 标题 / 副题 / 署名 / 摘要 / 关键词) ----
% 用法: \reporthead{kicker-left}{kicker-right}{标题}{副题}{署名行}
\newcommand{\reporthead}[5]{%
  \begin{center}
  {\small\sffamily\bfseries\color{accent} #1 \hfill #2}\\[2pt]
  {\color{black}\rule{\linewidth}{1.2pt}}\\[14pt]
  {\LARGE\sffamily\bfseries #3\par}\vspace{8pt}
  {\large\color{faint} #4\par}\vspace{10pt}
  {\small\sffamily\color{faint} #5\par}\vspace{4pt}
  {\color{rule}\rule{\linewidth}{0.4pt}}
  \end{center}\vspace{6pt}
}
% 摘要块: \begin{cnabstract} ... \end{cnabstract}
\newenvironment{cnabstract}{%
  \begin{quote}\small\color{faint}\sffamily\bfseries\color{black} 摘要 \quad\normalfont\rmfamily\color{faint!80!black}%
}{\end{quote}\vspace{2pt}}
% 重点结论框: \begin{quotebox} ... \end{quotebox}
\newenvironment{quotebox}{%
  \begin{quote}\small
  {\color{rule}\rule{\linewidth}{0.8pt}}\par\vspace{2pt}%
}{\par\vspace{2pt}{\color{rule}\rule{\linewidth}{0.8pt}}\end{quote}\vspace{2pt}}

\newcommand{\keywords}[1]{\begin{quote}\small\color{faint}{\bfseries\sffamily\color{black} 关键词:\ } #1\end{quote}}

% 表内彩色标记
\newcommand{\upgood}[1]{\textcolor{good}{\bfseries #1}}
\newcommand{\downbad}[1]{\textcolor{warn}{\bfseries #1}}
\newcommand{\dimtext}[1]{\textcolor{faint}{#1}}

\setlength{\parskip}{2pt}
\linespread{1.12}
\emergencystretch=3em  % CJK 段落容忍度,消除 overfull

% 编译: tectonic(XeTeX 引擎),需要系统安装 Noto Serif/Sans CJK SC 字体
\usepackage[a4paper,top=2.7cm,bottom=2.7cm,left=2.5cm,right=2.5cm,headheight=15pt]{geometry}
\usepackage{fontspec}
% 正文:拉丁衬线 + 思源宋体;标题:拉丁无衬线 + 思源黑体
% 按文件名从 TeX 字体树加载(tectonic bundle 自带 TeX Gyre;系统字体需 Noto CJK SC)
\setmainfont{texgyretermes}[Extension=.otf,UprightFont=*-regular,BoldFont=*-bold,ItalicFont=*-italic,BoldItalicFont=*-bolditalic]
\setsansfont{texgyreheros}[Extension=.otf,UprightFont=*-regular,BoldFont=*-bold,ItalicFont=*-italic,BoldItalicFont=*-bolditalic]
\setmonofont{texgyrecursor}[Extension=.otf,UprightFont=*-regular,BoldFont=*-bold,Scale=0.88]
\setCJKmainfont{Noto Serif CJK SC}[BoldFont={Noto Serif CJK SC Bold}]
\setCJKsansfont{Noto Sans CJK SC}[BoldFont={Noto Sans CJK SC Bold}]
\setCJKmonofont{Noto Sans Mono CJK SC}
% 带圈数字 ①-⑳ 等划入 CJK 类,用 CJK 字体渲染(xeCJK 默认按拉丁处理)
\xeCJKDeclareCharClass{CJK}{"2460 -> "24FF}

\usepackage{amsmath,amssymb,bm}
\usepackage{booktabs}
\usepackage{graphicx}
\usepackage{xcolor}
\usepackage{titlesec}
\usepackage{fancyhdr}
\usepackage{enumitem}
\usepackage{caption}
\usepackage{listings}

% ---- 配色(主色可改) ----
\definecolor{accent}{HTML}{0E5FA8}     % 强调蓝(链接/节号)
\definecolor{good}{HTML}{0E7A4D}       % 正向结果
\definecolor{warn}{HTML}{B45309}       % 负向/警示
\definecolor{faint}{HTML}{6B7889}      % 弱化文字
\definecolor{codebg}{HTML}{F0F2F5}
\definecolor{rule}{HTML}{C9CFD8}

\usepackage[colorlinks=true,linkcolor=accent,citecolor=accent,urlcolor=accent]{hyperref}

% ---- 章节标题:无衬线粗体,节号着色 ----
\titleformat{\section}{\Large\bfseries\sffamily}{\textcolor{accent}{\thesection}}{0.6em}{}
\titleformat{\subsection}{\large\bfseries\sffamily}{\textcolor{accent}{\thesubsection}}{0.5em}{}
\titleformat{\subsubsection}{\normalsize\bfseries\sffamily}{\textcolor{accent}{\thesubsubsection}}{0.5em}{}
\titlespacing{\section}{0pt}{1.6em}{0.7em}
\titlespacing{\subsection}{0pt}{1.2em}{0.5em}

% ---- 页眉页脚(Cosmos 式:左标题右日期,下细线) ----
\newcommand{\runningtitle}{}  % 主文档用 \renewcommand 设置
\newcommand{\reportdate}{}
\newcommand{\reportfooter}{}
\fancypagestyle{cnreport}{%
  \fancyhf{}
  \fancyhead[L]{\small\sffamily \runningtitle}
  \fancyhead[R]{\small\sffamily \reportdate}
  \fancyfoot[L]{\small\color{faint}\sffamily \reportfooter}
  \fancyfoot[R]{\small\sffamily \thepage}
  \renewcommand{\headrulewidth}{0.4pt}
  \renewcommand{\footrulewidth}{0pt}
}
\pagestyle{cnreport}

% ---- 题注与列表 ----
\captionsetup{font=small,labelfont={bf,sf},skip=6pt}
\setlist{itemsep=2pt,topsep=4pt,parsep=0pt}

% ---- 代码 ----
\lstset{
  basicstyle=\ttfamily\small,
  backgroundcolor=\color{codebg},
  frame=single, framerule=0.3pt, rulecolor=\color{rule},
  xleftmargin=6pt, xrightmargin=6pt,
  breaklines=true, columns=fullflexible, keepspaces=true,
  aboveskip=8pt, belowskip=8pt,
}

% ---- 报告头(kicker / 标题 / 副题 / 署名 / 摘要 / 关键词) ----
% 用法: \reporthead{kicker-left}{kicker-right}{标题}{副题}{署名行}
\newcommand{\reporthead}[5]{%
  \begin{center}
  {\small\sffamily\bfseries\color{accent} #1 \hfill #2}\\[2pt]
  {\color{black}\rule{\linewidth}{1.2pt}}\\[14pt]
  {\LARGE\sffamily\bfseries #3\par}\vspace{8pt}
  {\large\color{faint} #4\par}\vspace{10pt}
  {\small\sffamily\color{faint} #5\par}\vspace{4pt}
  {\color{rule}\rule{\linewidth}{0.4pt}}
  \end{center}\vspace{6pt}
}
% 摘要块: \begin{cnabstract} ... \end{cnabstract}
\newenvironment{cnabstract}{%
  \begin{quote}\small\color{faint}\sffamily\bfseries\color{black} 摘要 \quad\normalfont\rmfamily\color{faint!80!black}%
}{\end{quote}\vspace{2pt}}
% 重点结论框: \begin{quotebox} ... \end{quotebox}
\newenvironment{quotebox}{%
  \begin{quote}\small
  {\color{rule}\rule{\linewidth}{0.8pt}}\par\vspace{2pt}%
}{\par\vspace{2pt}{\color{rule}\rule{\linewidth}{0.8pt}}\end{quote}\vspace{2pt}}

\newcommand{\keywords}[1]{\begin{quote}\small\color{faint}{\bfseries\sffamily\color{black} 关键词:\ } #1\end{quote}}

% 表内彩色标记
\newcommand{\upgood}[1]{\textcolor{good}{\bfseries #1}}
\newcommand{\downbad}[1]{\textcolor{warn}{\bfseries #1}}
\newcommand{\dimtext}[1]{\textcolor{faint}{#1}}

\setlength{\parskip}{2pt}
\linespread{1.12}
\emergencystretch=3em  % CJK 段落容忍度,消除 overfull

% 编译: tectonic(XeTeX 引擎),需要系统安装 Noto Serif/Sans CJK SC 字体
\usepackage[a4paper,top=2.7cm,bottom=2.7cm,left=2.5cm,right=2.5cm,headheight=15pt]{geometry}
\usepackage{fontspec}
% 正文:拉丁衬线 + 思源宋体;标题:拉丁无衬线 + 思源黑体
% 按文件名从 TeX 字体树加载(tectonic bundle 自带 TeX Gyre;系统字体需 Noto CJK SC)
\setmainfont{texgyretermes}[Extension=.otf,UprightFont=*-regular,BoldFont=*-bold,ItalicFont=*-italic,BoldItalicFont=*-bolditalic]
\setsansfont{texgyreheros}[Extension=.otf,UprightFont=*-regular,BoldFont=*-bold,ItalicFont=*-italic,BoldItalicFont=*-bolditalic]
\setmonofont{texgyrecursor}[Extension=.otf,UprightFont=*-regular,BoldFont=*-bold,Scale=0.88]
\setCJKmainfont{Noto Serif CJK SC}[BoldFont={Noto Serif CJK SC Bold}]
\setCJKsansfont{Noto Sans CJK SC}[BoldFont={Noto Sans CJK SC Bold}]
\setCJKmonofont{Noto Sans Mono CJK SC}
% 带圈数字 ①-⑳ 等划入 CJK 类,用 CJK 字体渲染(xeCJK 默认按拉丁处理)
\xeCJKDeclareCharClass{CJK}{"2460 -> "24FF}

\usepackage{amsmath,amssymb,bm}
\usepackage{booktabs}
\usepackage{graphicx}
\usepackage{xcolor}
\usepackage{titlesec}
\usepackage{fancyhdr}
\usepackage{enumitem}
\usepackage{caption}
\usepackage{listings}

% ---- 配色(主色可改) ----
\definecolor{accent}{HTML}{0E5FA8}     % 强调蓝(链接/节号)
\definecolor{good}{HTML}{0E7A4D}       % 正向结果
\definecolor{warn}{HTML}{B45309}       % 负向/警示
\definecolor{faint}{HTML}{6B7889}      % 弱化文字
\definecolor{codebg}{HTML}{F0F2F5}
\definecolor{rule}{HTML}{C9CFD8}

\usepackage[colorlinks=true,linkcolor=accent,citecolor=accent,urlcolor=accent]{hyperref}

% ---- 章节标题:无衬线粗体,节号着色 ----
\titleformat{\section}{\Large\bfseries\sffamily}{\textcolor{accent}{\thesection}}{0.6em}{}
\titleformat{\subsection}{\large\bfseries\sffamily}{\textcolor{accent}{\thesubsection}}{0.5em}{}
\titleformat{\subsubsection}{\normalsize\bfseries\sffamily}{\textcolor{accent}{\thesubsubsection}}{0.5em}{}
\titlespacing{\section}{0pt}{1.6em}{0.7em}
\titlespacing{\subsection}{0pt}{1.2em}{0.5em}

% ---- 页眉页脚(Cosmos 式:左标题右日期,下细线) ----
\newcommand{\runningtitle}{}  % 主文档用 \renewcommand 设置
\newcommand{\reportdate}{}
\newcommand{\reportfooter}{}
\fancypagestyle{cnreport}{%
  \fancyhf{}
  \fancyhead[L]{\small\sffamily \runningtitle}
  \fancyhead[R]{\small\sffamily \reportdate}
  \fancyfoot[L]{\small\color{faint}\sffamily \reportfooter}
  \fancyfoot[R]{\small\sffamily \thepage}
  \renewcommand{\headrulewidth}{0.4pt}
  \renewcommand{\footrulewidth}{0pt}
}
\pagestyle{cnreport}

% ---- 题注与列表 ----
\captionsetup{font=small,labelfont={bf,sf},skip=6pt}
\setlist{itemsep=2pt,topsep=4pt,parsep=0pt}

% ---- 代码 ----
\lstset{
  basicstyle=\ttfamily\small,
  backgroundcolor=\color{codebg},
  frame=single, framerule=0.3pt, rulecolor=\color{rule},
  xleftmargin=6pt, xrightmargin=6pt,
  breaklines=true, columns=fullflexible, keepspaces=true,
  aboveskip=8pt, belowskip=8pt,
}

% ---- 报告头(kicker / 标题 / 副题 / 署名 / 摘要 / 关键词) ----
% 用法: \reporthead{kicker-left}{kicker-right}{标题}{副题}{署名行}
\newcommand{\reporthead}[5]{%
  \begin{center}
  {\small\sffamily\bfseries\color{accent} #1 \hfill #2}\\[2pt]
  {\color{black}\rule{\linewidth}{1.2pt}}\\[14pt]
  {\LARGE\sffamily\bfseries #3\par}\vspace{8pt}
  {\large\color{faint} #4\par}\vspace{10pt}
  {\small\sffamily\color{faint} #5\par}\vspace{4pt}
  {\color{rule}\rule{\linewidth}{0.4pt}}
  \end{center}\vspace{6pt}
}
% 摘要块: \begin{cnabstract} ... \end{cnabstract}
\newenvironment{cnabstract}{%
  \begin{quote}\small\color{faint}\sffamily\bfseries\color{black} 摘要 \quad\normalfont\rmfamily\color{faint!80!black}%
}{\end{quote}\vspace{2pt}}
% 重点结论框: \begin{quotebox} ... \end{quotebox}
\newenvironment{quotebox}{%
  \begin{quote}\small
  {\color{rule}\rule{\linewidth}{0.8pt}}\par\vspace{2pt}%
}{\par\vspace{2pt}{\color{rule}\rule{\linewidth}{0.8pt}}\end{quote}\vspace{2pt}}

\newcommand{\keywords}[1]{\begin{quote}\small\color{faint}{\bfseries\sffamily\color{black} 关键词:\ } #1\end{quote}}

% 表内彩色标记
\newcommand{\upgood}[1]{\textcolor{good}{\bfseries #1}}
\newcommand{\downbad}[1]{\textcolor{warn}{\bfseries #1}}
\newcommand{\dimtext}[1]{\textcolor{faint}{#1}}

\setlength{\parskip}{2pt}
\linespread{1.12}
\emergencystretch=3em  % CJK 段落容忍度,消除 overfull
